% !PRINTED 643
\weberpage{643}
If one sets therein
\[
\Delta = \pm\,q^{\,q^{\varkappa-1}[\varkappa(q-1)-1]}
\]
it follows that
\weq{7}{
\begin{array}{rl}
\Delta_1 = q^{\varkappa\, q^{\varkappa-1}\frac{q-1}{2} - \frac{q^{\varkappa-1}+1}{2}} & \text{for odd } q\\
= 2^{(\varkappa-1)2^{\varkappa-2}-1} & \text{for } q=2,
\end{array}
}
so that $\Delta_1$ is always a power of $q$.

\begin{center}
\S.\,177.\\[4pt]
\gesperrt{The Prime Ideals in the Field $H_m$.}
\end{center}

We have now to investigate the question in what relation the prime ideals of the field $H_m$ stand to those of the field $\Omega_m$.

If $\mathfrak{P}$ denotes any prime ideal in $H_m$ of degree $f_1$, then $\mathfrak{P}$ must be a divisor of a natural prime $p$, and $\mathfrak{P}$ can therefore (by \S.\,171) be divisible by the second or a higher power of a prime factor $\mathfrak{p}$ in $\Omega_m$ only when $p=q$.

But if $\mathfrak{P}$ is a divisor of $q$, it must be a power of $\sigma=1-r$. We set, say,
\[
\mathfrak{P}=\sigma^\lambda.
\]
If we take the norm of this with respect to $\Omega_m$, which is the square of the norm with respect to $H_m$, it follows (by \S.\,169) that
\[
q^{2f_1}=q^\lambda,
\]
hence $f_1=\frac12\lambda$. Since $f_1$ is an integer, $\lambda$ must be at least $=2$. But $\lambda$ also cannot be greater than $2$, since $\sigma^2$ is associated with the integer $(1-r)(1-r^{-1})$ contained in $H_m$, and hence $\sigma^2$ must be divisible by $\mathfrak{P}$. Accordingly $f_1=1$, and we obtain the theorem:

1. \gesperrt{The prime $q$ is the $\frac12\mu$th power of a prime of the first degree existing in $H_m$.}

Now let $p$ be different from $q$, and $\mathfrak{p}$ a prime factor of $\mathfrak{P}$ in the field $\Omega_m$ of degree $f$, which goes over into $\mathfrak{p}'$ under the substitution $(r,r^{-1})$. Then $\mathfrak{P}$ is divisible both by $\mathfrak{p}$ and by $\mathfrak{p}'$; on the other hand, $\mathfrak{p}\mathfrak{p}'$ is contained in $H_m$ and hence divisible by $\mathfrak{P}$, and we have to distinguish whether $\mathfrak{p}, \mathfrak{p}'$ are different or not.

If $\mathfrak{p}$ is different from $\mathfrak{p}'$, then $\mathfrak{P}$ is divisible by $\mathfrak{p}\mathfrak{p}'$, and it follows that $\mathfrak{P}=\mathfrak{p}\mathfrak{p}'$, and by formation of the norm $f_1=f$.
